% !TeX root = tesis.tex

A pesar de su importancia para la interoperabilidad entre blockchains, los \Bridges han demostrado ser uno de los componentes más vulnerables del ecosistema. En los últimos años se han producido numerosos ataques que resultaron en pérdidas económicas de gran magnitud. Se estima que, desde 2020, los ataques a \Bridges han ocasionado pérdidas por miles de millones de dólares, incluyendo incidentes individuales de cientos de millones \cite{belenkov2025sokreviewcrosschainBridge}.

Existen dos razones principales por las cuales estos sistemas constituyen objetivos particularmente atractivos para atacantes. En primer lugar, muchos diseños concentran grandes cantidades de activos en contratos custodios, lo que implica que una vulnerabilidad puede comprometer fondos significativos. En segundo lugar, los \Bridges presentan una superficie de ataque amplia, ya que involucran múltiples componentes: contratos inteligentes desplegados en dos blockchains distintas, así como infraestructura \Offchain encargada de transmitir eventos entre ellas. 

Diversas clases de vulnerabilidades han sido explotadas en sistemas de interoperabilidad. Una de ellas corresponde a errores en la verificación de pruebas criptográficas utilizadas para demostrar eventos entre cadenas. Un ejemplo representativo ocurrió en el ataque al \emph{Binance Bridge} en 2022, donde una falla en la verificación de Merkle Proofs permitió construir una prueba fraudulenta correspondiente a una transacción inexistente, y como consecuencia, el contrato en la cadena destino liberó activos que nunca habían sido bloqueados en la cadena de origen.

Otra categoría relevante corresponde a errores lógicos en los contratos inteligentes que implementan los componentes del \Bridge. El ataque al \emph{Nomad Bridge} en 2022 explotó una inicialización incorrecta del contrato que permitía aceptar mensajes inválidos como válidos. Una vez descubierto el exploit, múltiples usuarios replicaron la transacción maliciosa, lo que produjo pérdidas masivas en pocas horas.

También se han observado vulnerabilidades asociadas a actualizaciones del sistema o código heredado. En el caso del protocolo \emph{Qubit Finance}, inconsistencias entre funciones antiguas y nuevas del contrato permitieron emitir eventos de depósito sin que los activos correspondientes hubieran sido transferidos realmente, lo que llevó a la creación de tokens sin respaldo.

Otra fuente frecuente de incidentes es el compromiso de claves privadas utilizadas por operadores del \Bridge. Muchos sistemas dependen de validadores que firman mensajes para autorizar transferencias entre cadenas. Si un atacante obtiene suficientes claves privadas, puede generar mensajes fraudulentos que el sistema aceptará como válidos. Este tipo de ataque ocurrió, por ejemplo, en el \emph{Ronin Bridge}, donde los atacantes comprometieron múltiples validadores y autorizaron transferencias fraudulentas por un valor cercano a los 600 millones de dólares. 

La recurrencia de estos ataques ha motivado el desarrollo de arquitecturas de interoperabilidad más robustas, ya que las vulnerabilidades explotadas no corresponden a debilidades criptográficas fundamentales, sino a errores de implementación o supuestos de confianza incorrectos. En particular, un \zkBridge reduce la dependencia en operadores privilegiados y disminuyen la superficie de ataque del sistema. En la siguiente sección analizaremos esta arquitectura con mayor detalle.
